\chapter{Asterixis (Flapping Tremor)}

\section*{Definition}
Asterixis is a sudden, brief, irregular loss of sustained muscle contraction, a form of negative myoclonus. It is often seen as brief lapses of posture when the arms and wrists are held out, but it can affect other muscle groups and is not specific to liver disease. A clinician usually elicits it with the arms extended, wrists bent upward, and fingers spread.

\section*{When to See a Doctor}

\begin{itemize}
 \item New-onset asterixis, especially with liver, kidney, lung, or neurologic disease
 \item Asterixis with altered mental status, confusion, or disorientation
 \item Associated jaundice, ascites, or gastrointestinal bleeding
 \item Asterixis with fever, headache, or focal neurological deficit
\end{itemize}

\section*{How It Develops}
Asterixis reflects intermittent failure of brain networks that maintain posture and sustained muscle contraction. It most often signals a toxic or metabolic disturbance, but medications, respiratory failure, and focal brain lesions can also cause it. The exact pathway varies with the cause and should not be inferred from the sign alone.

\section*{Causes}
\begin{itemize}
 \item \textbf{Hepatic encephalopathy:} Most common cause; seen in acute or chronic liver failure
 \item \textbf{Uremia:} Chronic kidney disease, end-stage renal disease
 \item \textbf{Respiratory failure:} Hypercapnia (CO$_2$ retention), hypoxemia
 \item \textbf{Electrolyte disturbances:} Hyponatremia, hypokalemia, hypomagnesemia, hypocalcemia
 \item \textbf{Drug-induced:} Sedatives (benzodiazepines, barbiturates), anticonvulsants (phenytoin, valproate), lithium, levetiracetam
 \item \textbf{Structural brain lesions:} Thalamic stroke, midbrain hemorrhage, bilateral basal ganglia lesions
 \item \textbf{Other metabolic:} Wilson disease, hypoglycemia, Wernicke encephalopathy
\end{itemize}

\section*{Related Symptoms}
\begin{itemize}
 \item Bilaterally synchronous wrist flapping on sustained dorsiflexion
 \item May be elicited at the ankle (dorsiflexion) or by tongue protrusion
 \item Altered mental status, confusion, disorientation
 \item Fetor hepaticus (musty breath odor in liver failure)
 \item Jaundice, ascites, spider angiomata in hepatic disease
 \item Myoclonus, tremor, or other movement disorders
\end{itemize}
\textbf{Important:} Asterixis is a sign of metabolic encephalopathy, not a primary movement disorder; the underlying cause (hepatic, uremic, respiratory, or metabolic) must be identified and corrected.

\section*{Medical Evaluation}
\begin{itemize}
 \item Clinical examination, including posture, alertness, neurologic function, breathing, and medication or substance review
 \item Targeted laboratory evaluation may include glucose, electrolytes, calcium, kidney and liver tests, blood count, and blood gas testing when respiratory failure is possible. Ammonia may support an assessment of hepatic encephalopathy but does not confirm or exclude it by itself.
 \item Neuroimaging for focal neurologic findings, trauma, severe headache, or an unexplained cause; EEG may be considered when seizures are possible
 \item Lumbar puncture if meningoencephalitis is suspected
\end{itemize}

\section*{Treatment Options}
\begin{itemize}
 \item Treatment of the underlying metabolic derangement
 \item Hepatic encephalopathy: treat precipitants and use clinician-directed lactulose (often titrated to 2--3 soft stools/day) and rifaximin when indicated; maintain adequate nutrition rather than routinely restricting protein.
 \item Uremic encephalopathy: Dialysis, optimization of renal function
 \item Respiratory failure: Correction of hypercapnia and hypoxemia
 \item Electrolyte repletion as indicated
 \item Review and clinician-guided adjustment of causative medications; do not abruptly stop sedatives or antiseizure medicines without advice
 \item Liver transplantation for end-stage liver disease with refractory encephalopathy
\end{itemize}

\section*{Self-Care and Home Management}
\begin{itemize}
 \item Asterixis itself is not managed at home; it indicates an underlying metabolic emergency
 \item People with known liver disease should avoid unapproved sedatives and seek advice before changing medicines or diet; routine protein restriction is not recommended.
 \item Adherence to lactulose or rifaximin therapy in hepatic encephalopathy
 \item Maintain hydration and electrolyte balance
 \item Fall precautions due to postural instability
\end{itemize}

\section*{References}
\begin{thebibliography}{99}
\bibitem{asterixis:ref1} Frederick RT. ``Current concepts in the pathophysiology and management of hepatic encephalopathy.'' \textit{Gastroenterol Hepatol}. 2011;7(9):627-632.
\bibitem{asterixis:ref2} Vilstrup H, Amodio P, et al. ``Hepatic encephalopathy in chronic liver disease: 2014 practice guideline by AASLD and EASL.'' \textit{Hepatology}. 2014;60(2):715-735.
\bibitem{asterixis:ref3} Bajaj JS. ``Review article: the modern management of hepatic encephalopathy.'' \textit{Aliment Pharmacol Ther}. 2010;31(5):537-547.
\bibitem{asterixis:ref4} Schiff ER, Maddrey WC, et al. \textit{Schiff\'s Diseases of the Liver}, 12th ed. Wiley-Blackwell, 2017.
\bibitem{asterixis:ref5} UpToDate. ``Hepatic encephalopathy: clinical features and diagnosis.'' Wolters Kluwer, 2023.
\bibitem{asterixis:ref6} American Association for the Study of Liver Diseases. Diagnosis, evaluation, and management of hepatic encephalopathy in chronic liver disease. Updated guidance.
\bibitem{asterixis:ref7} Merck Manual Professional Edition. Myoclonus and asterixis. Updated 2025.
\end{thebibliography}
